\title{Parameter-Efficient Orthogonal Finetuning via Factorization}

\begin{abstract}
The growing prevalence of large foundation models has made their from-scratch training cost-prohibitive, prompting a need for more resource-conscious strategies for downstream adaptation. This paper investigates a systematic finetuning method—Orthogonal Finetuning (OFT)—for task adaptation using such models. While OFT boasts strong generalization, its reliance on high-dimensional orthogonal matrices leads to considerable numbers of tunable parameters. To mitigate this, we analyze OFT from the standpoint of information transmission and extract several criteria necessary for achieving greater parameter efficiency. Drawing inspiration from the Cooley-Tukey fast Fourier transform’s effectiveness in compactly transmitting information, we introduce an orthogonal parameterization based on butterfly architectures. Integrating this structure into OFT yields a novel, parameter-efficient finetuning technique, referred to as Orthogonal Butterfly~(BOFT). BOFT generalizes OFT as a specific instance, presenting a broader orthogonal finetuning schema. To validate our approach, we conduct comprehensive experiments on finetuning large vision transformers, large language models, and text-to-image diffusion models for a variety of vision and language downstream tasks.
\end{abstract}

\section{Introduction}

Recent developments, such as ChatGPT~\cite{brown2020language,chen2021evaluating} and Stable Diffusion~\cite{rombach2022high}, underscore the exceptional generalization capacity of large foundation models. This leap in capability has been accompanied by exponential growth in parameter counts 

(\emph{e.g.}, GPT-3 comprises approximately 175B parameters), which poses significant barriers for training these models anew. Thus, meaningful progress in the area relies on our ability to tailor such models to specific tasks without fully retraining them. Effectively, this requires parameter-efficient adaptation to downstream tasks. There are chiefly three avenues for efficient model adaptation: \emph{model finetuning}~\cite{hulora2022,zaken2022bitfit,guo2020parameter,raffel2020exploring,qiu2023controlling,zhang2022adaptive,chavan2023one}, wherein the model’s architecture remains static but only a portion of its parameters are updated; \emph{adapter tuning}~\cite{rebuffi2017learning,houlsby2019parameter,pfeiffer2020adapterfusion,lin2020exploring,he2021towards}, which augments the model with additional tunable components; and \emph{prompt tuning}~\cite{li2021prefix,lester2021power}, which appends learnable prefix tokens to the model’s inputs. Among these strategies, model finetuning stands out for its straightforwardness and effectiveness, with the added advantage of not increasing inference time.

Model finetuning fundamentally aims to maintain a high degree of similarity between the pretrained and the finetuned models, as quantified by specific metrics, in order to retain the pretrained knowledge. Standard finetuning practices, for example, commonly use a small learning rate as a heuristic, thereby limiting the extent to which the finetuned model deviates from the original. Recognizing the structured composition of weight matrices, a more systematic strategy seeks to preserve relational attributes within these matrices—specifically, the pairwise angular relationships among neurons. This rationale underpins the Orthogonal Finetuning~(OFT) framework~\cite{qiu2023controlling}, wherein all neurons within a layer are subjected to an identical orthogonal transformation, ensuring the pairwise angular configuration is mathematically conserved throughout finetuning. While OFT has yielded favorable results in terms of both generalization and optimization when finetuning text-to-image diffusion models~\cite{qiu2023controlling}, one significant drawback lies in the sheer number of parameters required to represent large orthogonal matrices in high dimensions. To mitigate this issue, OFT adopts a block-diagonal design to lower parameter count. Nevertheless, this improvement in parameter efficiency is offset by new challenges: the imposed block-diagonal structure enforces a rigid sparsity pattern and compartmentalizes orthogonal transformations across blocks, potentially introducing unwanted inductive biases (e.g., diminishing expressivity and limiting the ability to emulate conventional linear mappings). Moreover, the challenge of identifying optimal sparsity configurations within this framework remains unresolved.

A central challenge in addressing this issue is constructing a dense orthogonal matrix in a manner that remains efficient with respect to the number of parameters. Although at first it might appear that this is unattainable—since a fully dense, $d$-dimensional orthogonal matrix necessitates $\mathcal{O}(d^2)$ parameters—we introduce a novel approach: expressing the dense orthogonal matrix as a product of multiple sparse matrices in a factorized form. This strategy is predicated on the realization that it is possible to lower the parameter count by exchanging increased computational complexity for storage efficiency. Representing the orthogonal matrix as a composition of sparse matrices means the required matrix multiplications must be repeatedly carried out during each training step. To provide additional context for this factorization, we frame the creation of a dense orthogonal matrix as analogous to the process of transmitting information across a network. Specifically, this perspective recasts the construction of a dense orthogonal matrix via sparse matrix multiplication as an information dissemination process over a graph with grid-like connectivity. Adopting this viewpoint enables us to conceptualize numerous methods for sparse matrix factorization that carefully control parameter counts, while still retaining the capacity to realize dense matrix structures.

Within this information transmission paradigm, parameter efficiency is realized by drawing on the butterfly structures found in the Cooley-Tukey fast Fourier transform algorithm~\cite{cooley1965algorithm}, which efficiently facilitate the interchange of information among different nodes~\cite{klasing1994broadcasting}. In particular, the butterfly graph inherent to the Cooley-Tukey algorithm naturally leads to a form of sparse matrix factorization, referred to as butterfly factorization. Utilizing butterfly factorization allows the construction of a $d$-dimensional dense matrix through the product of $\mathcal{O}(\log d)$ sparse matrices, where each matrix contains only $\mathcal{O}(d)$ non-zero elements. By ensuring that every component sparse matrix remains orthogonal, we obtain a highly efficient orthogonal parameterization comprising just $\mathcal{O}(d\log d)$ parameters—a marked improvement over the naive quadratic scaling. Leveraging this parameter-efficient orthogonalization, we introduce a new finetuning technique—Orthogonal Butterfly~(BOFT)—which encompasses OFT as a specific instance and generalizes the concept of orthogonal finetuning. Both block-diagonal and butterfly structures share the feature of sparsity, thus embedding sparse patterns in orthogonal matrices and decreasing the number of parameters that must be learned. Our methodology stands in contrast with the low-rank decomposition technique employed by LoRA~\cite{hulora2022}; indeed, both low-rank and sparse matrices constitute significant classes within structured matrix theory~\cite{candes2011robust} and represent powerful approaches to parameter efficiency.

In contrast to OFT's block-diagonal parameterization, which balances model capacity and regularization, BOFT leverages the butterfly structure to enable a more seamless transition between elements from the full orthogonal group (offering greater expressivity) and the identity matrix (providing regularization). This design allows us to focus on a more refined subset of the orthogonal group suitable for specific downstream applications. Given the butterfly structure’s prominent role in accelerating numerous linear operations, including the discrete Fourier and discrete cosine transforms, we propose that this architectural choice not only enhances parameter efficiency, but also embeds an implicit inductive bias that can foster improved generalization and help mitigate overfitting. Our perspective is motivated by the butterfly structure’s capability to represent a wide array of classical linear operators—capabilities not achievable with the block-diagonal approach utilized in OFT. Our primary contributions can be summarized as follows:

\begin{itemize}[leftmargin=*,nosep]
    \item We examine the challenge of parameter-efficient orthogonal finetuning through a novel lens of information transmission, recasting the task of designing compact dense orthogonal matrices as an information flow problem over a grid-structured graph.
    \item Drawing inspiration from the butterfly networks underpinning the Cooley-Tukey algorithm, we introduce Orthogonal Butterfly, an orthogonal finetuning technique that prioritizes parameter efficiency within the information transmission paradigm.
    \item We present theoretical analyses elucidating how BOFT drastically decreases the number of trainable parameters while maintaining strong expressive capacity and training robustness. The matrix decomposition inherent in BOFT also yields a useful framework for interpolating model weights (refer to Figure~\ref{fig:boft_interpolation}).
    \item We pioneer the use of orthogonal finetuning across diverse domains beyond the original controllable text-to-image generation task~\cite{qiu2023controlling}, demonstrating its viability as a general-purpose model adaptation technique. Specifically, we evaluate BOFT on a spectrum of applications, spanning computer vision to natural language processing, where it substantially surpasses existing state-of-the-art techniques, substantiating its high parameter efficiency and capacity for generalization.
\end{itemize}

\section{Related Work}

\textbf{Parameter-efficient finetuning (PEFT)}. As the scale and capabilities of foundation models (\emph{e.g.}, \cite{brown2020language,radford2021learning,rombach2022high,kirillov2023segment}) continue to grow, efficiently adapting these models to downstream applications with minimal parameter overhead has become a focal research area~\cite{treviso2023efficient,ding2023parameter,lialin2023scaling}. Numerous PEFT techniques have been introduced~\cite{houlsby2019parameter,aghajanyan2020intrinsic,hulora2022,edalati2022krona,wang2022adamix,gheini2021cross,zaken2022bitfit,guo2020parameter,sung2021training,ansell2022composable,lester2021power,li2021prefix,vu2022spot,he2021towards,mao2021unipelt,karimi2021compacter,liu2022few,sung2022lst,chen2022parameter,jia2022visual,chen2022adaptformer,zhang2022neural,jie2023fact,lian2022scaling,luo2023towards}; among these, reparameterization approaches~\cite{aghajanyan2020intrinsic,hulora2022,edalati2022krona} are most aligned with our research direction. A prominent example, LoRA~\cite{hulora2022}, enhances pretrained weight matrices by incorporating products of two low-rank matrices, demonstrating strong results for natural language processing tasks. However, LoRA enforces a constant rank across all layers; as a response, approaches such as \cite{zhang2022adaptive,valipour2022dylora,zhang2023increlora} have proposed dynamically adjusting the rank per layer to better distribute the parameter budget. The appeal of this straightforward low-rank decomposition is evidenced by its widespread adoption~\cite{dettmers2023qlora,zi2023delta,chavan2023one}. Taking inspiration from the way hyperspherical energy provides insight into model generalization~\cite{liu2018learning,liu2021orthogonal}, \cite{qiu2023controlling} introduced orthogonal finetuning (OFT)—a contrasting yet robust approach designed for finetuning text-to-image diffusion models. In this framework, an orthogonal matrix is learned to linearly transform neurons within each layer, leading to improved generalization and more consistent training relative to LoRA. Nevertheless, despite these advantages, OFT typically requires more trainable parameters compared to LoRA, pointing to the need for greater parameter efficiency. Additionally, the potential for applying OFT more broadly—beyond the context of text-to-image diffusion model adaptation—remains unexplored. BOFT addresses this gap by leveraging butterfly factorization to make OFT more parameter efficient, thus unlocking orthogonal finetuning’s benefits for a wider array of adaptation tasks.

\textbf{Butterfly structures}. The radix-2 Cooley-Tukey algorithm achieves computational efficiency for the $N$-point discrete Fourier transform by iteratively decomposing it into pairs of $\frac{N}{2}$-point transforms, resulting in a characteristic butterfly pattern; this can be mathematically represented as a product of sparse matrices (collectively known as a butterfly matrix). Such butterfly matrices have been employed to orthogonally parametrize matrices for improved Gaussian elimination (eliminating the need for pivoting)~\cite{parker1995random}, to bolster the stability of training recurrent neural networks~\cite{jing2017tunable}, and for efficient kernel approximations~\cite{munkhoeva2018quadrature}. Learning rapid linear transforms with butterfly parameterizations has been explored by \cite{dao2019learning,dao2019kaleidoscope}, while \cite{chen2022pixelated,dao2022monarch} utilize them for optimizing neural network training efficiency. Applications extend further, including fast matrix-vector multiplication~\cite{michielssen1996multilevel,o2010algorithm}, compact data-sparse matrix representations~\cite{li2015butterfly}, and network communication protocols~\cite{klasing1994broadcasting,li2003linear,rajasingh2016transmission}. Distinct from these prior efforts, our work focuses specifically on employing butterfly structures to boost the parameter efficiency of OFT.

\section{An Information Transmission View on Orthogonal Finetuning}

We begin by reviewing key aspects of OFT. OFT adapts the pretrained weight matrix by expressing the adapted matrix as a product of a trainable orthogonal matrix and the fixed pretrained weight matrix. Unlike LoRA, which fine-tunes the weights by adding a low-rank matrix, OFT employs a multiplicative approach, using a trainable orthogonal matrix. LoRA attains parameter efficiency via a low-rank constraint, whereas the standard OFT achieves efficiency through a block-diagonal orthogonal structure~\cite{qiu2023controlling}; smaller diagonal block sizes lead to greater parameter efficiency. This distinction is illustrated in Figure~\ref{fig:comp_lora}. The underlying motivation for leveraging orthogonal transformations when tuning the weight matrix is to maintain the pairwise angular relationships between neurons~\cite{liu2018learning,liu2021orthogonal,qiu2023controlling}, thus better retaining the semantic knowledge acquired during pretraining. Specifically, OFT learns an orthogonal matrix $\bm{R}\in\mathbb{R}^{d\times d}$ for a given pretrained linear transformation $\bm{W}^0\in\mathbb{R}^{d\times n}$, changing the computation from $\bm{z}=(\bm{W}^0)^\top\bm{x}$ to $\bm{z}=(\bm{R}\bm{W}^0)^\top\bm{x}$, with $\bm{x}\in\mathbb{R}^{d}$ and $\bm{z}\in\mathbb{R}^{n}$ representing input and output vectors, respectively. To ensure zero initialization, $\bm{R}$ is initialized to the identity matrix. Preservation of orthogonality during training is maintained using Cayley parameterization, following~\cite{qiu2023controlling,liu2021orthogonal}: $\bm{R}=(\bm{I}+\bm{Q})(\bm{I}-\bm{Q})^{-1}$, where $\bm{Q}$ denotes a skew-symmetric matrix satisfying $\bm{Q}=-\bm{Q}^\top$. To improve parameter efficiency, the orthogonal matrix $\bm{R}$ is further parameterized as block-diagonal: $\text{diag}(\bm{R}_1,\bm{R}_2,\cdots,\bm{R}_r)$, where each block $\bm{R}_i\in\mathcal{R}^{b\times b}$ for all $i$, and $br=d$. However, while this block-diagonal structure enhances parameter efficiency, it imposes the assumption that a neuron's dimensions (i.e., a column in $\bm{W}^0$) are split into $r$ separate groups, each subject to a distinct orthogonal transformation, despite these partitions being determined only by index order and not by semantic or functional relevance. Although this approach is effective in practice, there is little justification for dividing neuron dimensions solely by indices, motivating the exploration of dense orthogonal matrices. This raises an important question: \emph{Is it possible to design a dense orthogonal matrix structure that retains parameter efficiency?}

In response to this issue, we introduce an approach where a dense orthogonal matrix is constructed by multiplying several sparse orthogonal matrices. To systematically investigate the sparsity configuration inherent in orthogonal matrix factorization, we reinterpret the creation of a dense orthogonal matrix within the OFT framework as an information transmission task. Concretely, forming a dense matrix $\bm{R}\in\mathbb{R}^{d\times d}$ as the product $\bm{R} = \bm{B}_m\bm{B}_{m-1}\cdots\bm{B}_1$ with $m$ square matrices can be analogized to transmitting information across a grid structured by $d$ rows and $m+1$ columns, as depicted in Figure~\ref{fig:info_trans}. The rationale for adopting this information-theoretic lens arises from the fact that a dense $d$-by-$d$ matrix can be viewed as specifying complete communication links between two sets of $d$ nodes. Within the matrix $\bm{R}$, a zero entry $R_{ij}$ signifies that information originating from node $j$ does not reach node $i$; conversely, a nonzero $R_{ij}$ permits such transmission. Consequently, the representation of $\bm{R}$ as a sequence of multiplications $\bm{B}_m\bm{B}_{m-1}\cdots\bm{B}_1$ naturally corresponds to staged information flow governed by the adjacency graphs associated with each $\bm{B}_i$. Transmission proceeds in accordance with the structure imposed by $\bm{B}_1$, culminating with $\bm{B}_n$. As illustrated in Figure~\ref{fig:info_trans}, for the decomposition $\bm{R} = \bm{B}_5\bm{B}_4\bm{B}_3\bm{B}_2\bm{B}_1$, both the associated sparsity patterns and the resulting graph are visualized. The graph representation effectively ‘unrolls’ the steps in the matrix product, interpreting each $\bm{B}_i$ as defining directed edges from nodes at level $i$ to nodes at level $i+1$. Specifically, entry $(j_1,j_2)$ of $\bm{B}_i$ is nonzero if and only if there exists a directed link from node $j_2$ in the $i$-th layer to node $j_1$ in the $(i+1)$-th layer. For $\bm{R} = \bm{B}_5\bm{B}_4\bm{B}_3\bm{B}_2\bm{B}_1$ to be fully dense, it is necessary that information from any node at the initial layer can reach any node at the sixth layer. Examining, for example, the product $\bm{R} = \bm{B}_4\bm{B}_3\bm{B}_2\bm{B}_1$ (which entails only layers one through five), one observes that data originating from node 1 fails to reach node 3, revealing a zero at position $(3,1)$ in the associated product’s sparsity pattern. Similar observations apply for products of fewer factors, such as $\bm{R} = \bm{B}_3\bm{B}_2\bm{B}_1$ and $\bm{R} = \bm{B}_2\bm{B}_1$, maintaining the correspondence between the graphs and the sparsity of the resultant matrices. Broadly, in order to achieve a dense matrix $\bm{R}\in\mathbb{R}^{d\times d}$, the factorization matrices $\bm{B}_m,\dots,\bm{B}_1$ must collectively form directed edges on a $d\times (m+1)$ node grid, with each edge spanning only adjacent columns. This guarantees that each source node at the first layer can transmit to every destination node at the $(m+1)$-th layer.

Adopting the perspective of information transmission provides further insight into the sparse patterns present in the factorization matrices utilized in OFT. As depicted in Figure~\ref{fig:blk_diag}, the original OFT employs a block-diagonal configuration for $\bm{R}$. While this approach leads to a reduced count of learnable parameters, it prevents the resulting matrix $\bm{R}$ from achieving full density. The primary objective, then, is to form a dense orthogonal matrix by combining $m$ sparse orthogonal matrices, aiming to do so while minimizing the number of effective trainable parameters.

Through the lens of information transmission, two principal requirements naturally arise: \emph{(i)} \textbf{dense connectivity}, wherein each node at the input layer maintains at least one path to every node at the output layer, and \emph{(ii)} \textbf{minimal free edges}, which dictates that the overall edge count should be as low as possible without violating the orthogonality condition. Importantly, orthogonality imposes nuanced restrictions on edge formation between subsequent levels. To illustrate, for every matrix $\bm{B}_i$ to have full rank—a prerequisite for orthogonality—it is necessary to establish $d$ edges that correspond to a bijective mapping between nodes in level $i$ and nodes in level $(i=1)$. Thus, at least $d$ edges must link any two consecutive levels (\emph{e.g.}, this is 4 in the instance shown in Figure~\ref{fig:info_trans}). Since these edges are dictated by the orthogonality condition and are non-trainable (for instance, a $d\times d$ orthogonal matrix with $d$ non-zero elements can only assign these to $\pm 1$), they should not be factored into the count of trainable edges. 

Owing to the lower parameter requirements of orthogonal matrices compared to unrestricted ones, the orthogonality property introduces additional dependencies among the edges. For example, in the context of a $2\times 2$ orthogonal matrix, assigning a zero to entry $(1,1)$ automatically necessitates a zero at $(2,2)$; that is, the presence or absence of a single edge can directly influence others. Consequently, orthogonality may compel the addition or removal of certain connections for any particular configuration. By examining the pattern of non-zero entries within the constructed orthogonal matrix, the information transmission view becomes especially informative for OFT, since only the existence of non-zero, trainable elements in $\bm{R}$ is significant, while their exact magnitudes are irrelevant.

A straightforward approach to creating dense connections between two layers requires $\mathcal{O}(d^2)$ edges (essentially, forming a fully dense orthogonal matrix), which specifically amounts to $d^2 - d$ edges; for instance, with $d=4$, there are 12 edges. Figure~\ref{fig:info_trans} illustrates a practical example of matrix factorization using this idea, consuming only $10$ edges overall—fewer than a standard dense orthogonal matrix would require. This representation provides a foundation for analyzing OFT’s parameter efficiency from the perspective of graph theory, and it allows us to construct valid matrix factorizations with ease. Motivated by the structure of the Cooley-Tukey algorithm, particularly its use of butterfly graphs~\cite{cooley1965algorithm}, we observe that such a topology achieves full connectivity between $d$ input and $d$ output nodes using only $\mathcal{O}(d\log d)$ edges. To illustrate, while the configuration in Figure~\ref{fig:info_trans} achieves dense connections with $10$ edges, a butterfly network would achieve the same with just $8$ edges. In the following, we explore how utilizing the butterfly architecture can enhance parameter efficiency.

\section{Orthogonal Parameterization by Butterfly Factorization}

Originating from the Cooley-Tukey algorithm for efficient fast Fourier transforms, the butterfly structure enables local changes in the frequency space to result in global effects in the spatial domain—mirroring the situation in our information transmission scenario, where any node at the initial layer can communicate with all nodes in the terminal layer. Butterfly structures are also well-established as network topologies for optimized information dissemination~\cite{solihin2015fundamentals,li2011linear}. Letting $k\geq 2$ be a power of $2$, we formally define the \emph{butterfly factor} $\bm{B}^F(k)$ as follows:

\begin{equation}\label{eq:but_factor}
\begin{aligned}
    \bm{B}^F(k)=\begin{bmatrix}
    \text{diag}(\bm{d}_1) & \text{diag}(\bm{d}_2) \\
    \text{diag}(\bm{d}_3) & \text{diag}(\bm{d}_4)
    \end{bmatrix}\in\mathbb{R}^{k\times k},
\end{aligned}
\end{equation}

where each $\bm{d}_i\in\mathbb{R}^{\frac{k}{2}},\forall i$ denotes a vector. Building upon this, and assuming $d=2^N$, we recursively define the $d$-dimensional \emph{butterfly matrix} $\bm{B}(d)\in\mathbb{R}^{d\times d}$ by

\begin{equation}
\begin{aligned}
    \!\!\bm{B}(d)=\tilde{\bm{B}}(d,d)\!\cdot\!\begin{bmatrix}
    \bm{B}_1(\frac{d}{2})\!\!\!\!\! & \bm{0} \\
    \bm{0}\!\!\!\!\! &\bm{B}_2(\frac{d}{2})
    \end{bmatrix}= \tilde{\bm{B}}(d,d)\tilde{\bm{B}}(d,\frac{d}{2})\cdots\tilde{\bm{B}}(d,2),
\end{aligned}
\end{equation}

where $\bm{B}_1(\frac{d}{2})$ and $\bm{B}_2(\frac{d}{2})$ denote two butterfly matrices each of dimension $\frac{d}{2}$. We introduce the \emph{butterfly component} as $\thickmuskip=2mu \medmuskip=2mu \tilde{\bm{B}}(d,k)=\text{diag}(\bm{B}^F_1(k),\cdots,\bm{B}^F_{d/k}(k))$, representing a block-diagonal matrix of dimensions $d\times d$ whose constituent blocks are of size $k$. Here, every $\bm{B}^F_i(k),\forall i$ corresponds to butterfly factors as described in Equation~\ref{eq:but_factor}. Utilizing this structure, we can parameterize an orthogonal matrix via butterfly matrices. This is accomplished by constraining all multiplicative factors $\tilde{\bm{B}}(d,k),\forall k$ within the butterfly matrix $\bm{B}(d)$ to be orthogonal matrices. 

Consider specifically the block-diagonal matrix $\tilde{\bm{B}}(d,2)$ with blocks of size $2$. Orthogonality for $\tilde{\bm{B}}(d,2)$ can be readily assured by adopting the Cayley transform (or representing each block with a $2$-dimensional rotation), consistent with the methods in \cite{liu2021orthogonal,qiu2023controlling}. The sparsity structure inherent to each butterfly component reflects a permutation of the sparsity pattern of $\tilde{\bm{B}}(d,2)$; as such, all butterfly components lend themselves to straightforward orthogonal parameterization. This approach facilitates a computationally efficient representation of orthogonal matrices that is based on numerous $2\times 2$ orthogonal blocks. 

Extending the butterfly matrices along the lines of \cite{chen2022pixelated}, we define the block butterfly component $\tilde{\bm{B}}^b(d,k)$, where every entry in $\bm{d}_i,\forall i$ is now a $b\times b$ block. To ensure that the block butterfly component $\tilde{\bm{B}}^b(d,2)$ maintains orthogonality, each constituent $2b\times 2b$ matrix block is parameterized as orthogonal. The blockwise non-zero arrangement of remaining butterfly components $\tilde{\bm{B}}^b(d,k), k>2$, can be viewed as permutations of the pattern in $\tilde{\bm{B}}^b(d,2)$, thus making their orthogonal parameterization equally feasible. Altogether, the BOFT forward computation can then be formulated as

\begin{equation}
    \begin{aligned}\nonumber
    \bm{z}=\big(\bm{R}(m,b)\cdot\bm{W}^0\big)^\top\bm{x},~~~~\text{s.t.}~\left\{\bm{R}(m,b)=\prod_{i=1}^m \tilde{\bm{B}}^b_{(d,i)}~~\&~~\underbrace{\big(\tilde{\bm{B}}^b_{(d,j)}\big)^\top\tilde{\bm{B}}^b_{(d,j)}=\tilde{\bm{B}}^b_{(d,j)}\big(\tilde{\bm{B}}^b_{(d,j)}\big)^\top\!=\bm{I}_d}_{\forall j\in [1, m]}\right\},
    \end{aligned}
\end{equation}

Here, for clarity, we define $\tilde{\bm{B}}^b(d,2^{m - i+1})$ as $\tilde{\bm{B}}^b_{(d,i)}$, with $\bm{I}_d$ denoting the $d$-dimensional identity matrix. The matrix $\bm{R}(m,b)\in\mathbb{R}^{d\times d}$ is constructed as a product of $m$ orthogonal butterfly matrices. For notational convenience, we abbreviate BOFT with $\bm{R}(m,\frac{b}{2})$ as BOFT($m$,$b$), under the condition $b\geq 2$. When $\thickmuskip=2mu \medmuskip=2mu m=1$, BOFT($1$,$b$) simplifies to the block-diagonal OFT~\cite{qiu2023controlling} with a block size of $b$. Furthermore, BOFT($1$,$d$) reduces to the standard OFT~\cite{qiu2023controlling} that utilizes an unconstrained full orthogonal matrix. The configuration BOFT($\log \frac{2d}{b}$,$b$) employs the block butterfly matrix $\bm{B}^{\frac{b}{2}}(d)$ as $\bm{R}$, which results in a dense orthogonal matrix $\bm{R}$. In a general setting, BOFT($m$,$b$) involves $\frac{1}{2}(b-1)dm$ trainable parameters for adapting a linear layer of dimensions $d\times n$. In the scenario where the butterfly structure is chosen, i.e., $m=\log d, b=2$, BOFT only requires $\mathcal{O}(d\log d)$ parameters. By contrast, the original OFT with a full dense orthogonal matrix entails $\mathcal{O}(d^2)$ parameters, and the block-diagonal OFT, assuming $r$ blocks, uses $\mathcal{O}(bd)$. This indicates that to obtain a dense orthogonal matrix, OFT needs to set $b=d$ as the block size, whereas BOFT has the flexibility to select any value for $b$.

\textbf{Identity initialization for BOFT}. Standard finetuning practices begin with the pretrained weights unchanged to ensure that modifications remain close to the pretrained solution. For example, in LoRA, low-rank adaptation weights are initialized to zero. In the BOFT framework, every butterfly component is initialized as an identity matrix (that is, skew-symmetric matrices in the Cayley transform are set to zero).

\textbf{Multiplicative Dropout for BOFT}. In LoRA~\cite{hulora2022}, Dropout is utilized to regularize low-rank updates and mitigate overfitting. Whereas standard Dropout~\cite{srivastava2014dropout} can be seamlessly applied to LoRA’s additive structure, BOFT's multiplicative nature renders traditional Dropout inapplicable. To address this, we introduce a multiplicative Dropout mechanism for BOFT. Noting that the orthogonal matrix $\bm{R}(m,b)$ comprises $m$ butterfly modules, each of which can be decomposed into $2b\times 2b$ block-diagonal orthogonal matrices, the procedure randomly selects a proportion $p_1$ of butterfly components and $p_2$ of the diagonal blocks within each, substituting the selected blocks by identity matrices.

\section{Intriguing Insights and Discussions}

\textbf{Expressivity of BOFT}. The butterfly framework, combined with permutations, is capable of exactly reproducing various well-known fast linear transforms~\cite{dao2019learning,dao2019kaleidoscope} (\emph{e.g.}, fast Fourier transform, Hadamard transform). However, the capacity of our orthogonal butterfly matrix to approximate arbitrary orthogonal matrices remains an open question. To investigate this, we simulate the approximation of a randomly generated dense orthogonal matrix~\cite{anderson1987generation} of size $512\times 512$. Figure~\ref{fig:para_eff} presents results averaged over 10 distinct random seeds. The y-axis captures the approximation error, while the x-axis reports the number of effective trainable parameters. Each color-coded curve corresponds to BOFT instances with identical block size, where the leftmost datapoint represents the error from BOFT($1$,$b$) (\emph{i.e.}, standard block-diagonal OFT with block size $b$). Across configurations, BOFT tends to achieve greater parameter efficiency relative to OFT. For instance, BOFT($9$,$2$) achieves higher expressivity than BOFT($1$,$16$) while requiring significantly fewer parameters. Configurations employing smaller $b$ values with larger $m$ (the number of butterfly layers) are typically more efficient in terms of parameter utilization; for example, BOFT($6$,$4$) attains comparable approximation error to BOFT($2$,$16$) using far fewer parameters. Broadly, the butterfly matrix imposes a richer structural constraint on the orthogonal group than the block-diagonal arrangement, making BOFT formally more expressive than OFT for equivalent block sizes.

\begin{theorem}[Expressivity of BOFT]\label{thm:exp_boft}
    BOFT surpasses OFT in expressivity when using the same block size. To fully approximate all orthogonal matrices of size $d$, butterfly matrices can be sequentially composed as $\bm{B}_{d-1,1}(d)\bm{B}^\top_{d-1,2}(d)\cdots\bm{B}_{1,1}(d)\bm{B}^\top_{1,2}(d)$, where each $\bm{B}_{i,j}(d)$ is a butterfly matrix.
\end{theorem}

Theorem~\ref{thm:exp_boft} implies a straightforward generalization of BOFT. Specifically, the final orthogonal matrix can be generalized as $\thickmuskip=2mu \medmuskip=2mu \bm{R}^G(m_1,b_1,m_2,b_2,l)=\bm{R}_{l,1}(m_1,b_1)\bm{R}_{l,2}^T(m_2,b_2)\cdots\bm{R}_{1,1}(m_1,b_1)\bm{R}_{1,2}^T(m_2,b_2)$, where $\bm{R}_{i,j}^T(m,b)$ refers to the orthogonal matrices within BOFT. In the scenario where $m_1=m_2=\log d$, $b_1=b_2=2$, and $l=d-1$, $\bm{R}^G(m_1,b_1,m_2,b_2,l)$ spans the full orthogonal group. This hierarchical composition, also known as the kaleidoscope hierarchy~\cite{dao2019kaleidoscope}, is theoretically powerful. However, greater expressiveness does not necessarily translate into better finetuning outcomes; for example, models trained with full finetuning—which has maximal expressivity—may exhibit suboptimal performance. Thus, balancing the degree of expressiveness and structural regularity is crucial for effective model finetuning. BOFT enlarges the search space for finetuning parameters via its structural prior, thereby offering a mechanism to attain a more optimal compromise between expressivity and regularity.

\textbf{Spectral properties}. Orthogonal finetuning typically results in more favorable spectral characteristics compared to LoRA, primarily because it maintains the spectral norm of the original pretrained weight matrix $\bm{W}^0$ without alteration. This can be clarified using singular value decomposition, where $\bm{W}^0=\bm{U}\bm{\Sigma}\bm{V}^\top$ with $\bm{U}$ and $\bm{V}$ as orthogonal matrices and $\bm{\Sigma}$ as a diagonal matrix containing singular values. In both OFT and BOFT, an orthogonal transformation $\bm{R}$ is applied to the left, producing adapted weights of the form $\bm{R}\bm{U}\bm{\Sigma}\bm{V}^\top$. This operation preserves the maximal singular value—namely, the spectral norm—of $\bm{W}^0$. The conservation of the spectral norm is well-known to promote stability in optimization and improve generalization capabilities during training~\cite{miyato2018spectral,yoshida2017spectral}. Additional mathematical analysis related to these spectral properties is available in Appendix~\ref{app:sn_property}.

\textbf{Orthogonal finetuning as learning bilinear similarity}. The BOFT approach can be interpreted as learning a bilinear similarity, specifically $\bm{w}^0_i\bm{R}\bm{x}$, where $\bm{w}^0_i$ denotes the $i$-th column (neuron) from the weight matrix $\bm{W}^0$. Within this framework, BOFT optimizes a similarity matrix $\bm{R}$ that is subject to orthogonality constraints, thereby establishing a connection with concepts in distance metric learning~\cite{xing2002distance} and bilinear forms~\cite{roman2005advanced}.

\textbf{Inductive bias and generalization in BOFT}. Since in BOFT the parameter $\bm{R}(m,b)$ generally denotes a specific structured subset of the orthogonal group, this imposes a restriction on the function class and thereby introduces an inductive bias. We posit that the kind of structured bias inherent to butterfly factorization aids generalization, as it mirrors structural patterns observed in various canonical linear transforms~\cite{dao2019kaleidoscope}, including the discrete Fourier, discrete sine/cosine, and Hadamard transforms. Furthermore, the implementation of sparse matrix factorization within BOFT is likely to confer additional implicit inductive biases~\cite{gunasekar2017implicit,li2018algorithmic,liu2019neural}.

\textbf{Comparison to butterfly-based sparse training}. A substantial body of prior research~\cite{dao2019learning,dao2019kaleidoscope,chen2022pixelated,dao2022monarch} has investigated sparse neural network training utilizing butterfly parameterizations. These approaches predominantly involve substituting conventional weight matrices with butterfly-structured matrices, which are then learned during training from initialization. Notably, \cite{dao2022monarch} extends this line by finetuning pretrained models via projection onto modified butterfly matrices, subsequently updating only the projected components for task adaptation. In contrast, BOFT introduces a novel finetuning paradigm that applies transformations to the weights through matrices that are shared across layers.

\input{rebuttal-results}

\section{Concluding Remarks and Limitations}

In this work, we introduce Orthogonal Butterfly, a parameter-efficient approach to finetuning foundation models grounded in the butterfly matrix structure. Our central contribution lies in the realization that improved parameter efficiency can be achieved by expressing a dense orthogonal matrix as a product of multiple sparse orthogonal matrices. To systematically identify such matrix decompositions, we develop a graphical framework for modeling information flow, through which we discover that butterfly matrices are particularly well-suited for constructing sparse orthogonal factorizations. BOFT is empirically shown to deliver strong performance when finetuning a diverse set of large-scale models—including language, vision, and text-to-image generators—thereby establishing its versatility and effectiveness as a general-purpose finetuning technique.

Nevertheless, BOFT is not without drawbacks. The requirement to compute products of several orthogonal matrices during finetuning leads to moderately increased training runtime compared to OFT, and methods for alleviating this overhead remain an area for future investigation. On the positive side, after finetuning, the learned orthogonal transformation can be absorbed directly into the original model, ensuring no additional cost at inference time. Additionally, the question of whether the butterfly structure provides the optimal information transmission mechanism is still open. Our information transmission framework points toward parallels with research in computer networking, a field where network topologies are analyzed for communication efficiency. This interdisciplinary insight suggests that exploring alternative network topologies could yield even more efficient strategies for dense orthogonal matrix construction.

\end{document}